% Week 8 section of the synthesis exercise sheet, included by Exercises.tex.
\section{Week 8 --- No-Arbitrage Principles and the Fundamental Theorems of Asset Pricing}

Throughout, \(B_t=e^{rt}\) denotes the money-market account under a constant
short rate \(r\), and \(P_t(T)\) denotes the time-\(t\) price of a
zero-coupon bond paying \(1\) at \(T\). In one-period examples, \(R\) denotes
the bond's gross return over the period.

\subsection*{Part I --- Dominance and the Law of One Price}

\begin{exercise}[Easy --- Turning a paid-to-enter trade into a free-and-sometimes-profitable one]
A one-period market has two states \(\omega_1,\omega_2\) and a bond with
price \(S_0^0=1\) and payoff \(R=1.08\) in both states. A strategy
\(\theta\) has
\[
V_0(\theta)=-4,
\qquad
V_T(\theta,\omega_1)=2,
\qquad
V_T(\theta,\omega_2)=0.
\]
\begin{exparts}
  \item Check that \(\theta\) is a type-(ii) arbitrage in the sense used in
  lecture.
  \item Augment \(\theta\) with the correct number of units of the bond so
  that the new strategy costs exactly \(0\) today. Compute its terminal
  value in each state and confirm it is a type-(i) arbitrage.
  \item Explain, in one sentence and without redoing the arithmetic, why the
  same construction works for \emph{any} strategy with \(V_0(\theta)<0\) and
  \(V_T(\theta,\cdot)\geq0\).
\end{exparts}
\end{exercise}


\begin{exercise}[Easy to medium --- Marking an existing forward to market]
A forward contract on a non-dividend-paying stock was entered at some past
date with a fixed delivery price \(K\) --- not necessarily the fair forward
price at inception. The contract still pays \(S_T-K\) at maturity \(T\).
\begin{exparts}
  \item Build the static replicating portfolio for this payoff at a current
  time \(t<T\), and use the law of one price to find its current value
  \(V_t\) in terms of \(S_t\), \(K\), \(r\), and \(T-t\).
  \item Check that your formula gives \(V_t=0\) when \(K\) equals the fair
  forward price \(K_t=S_te^{r(T-t)}\) quoted for a \emph{new} contract entered
  at \(t\).
  \item \emph{Numerical.} A long forward was entered when \(S_{t_0}=100\),
  \(r=4\%\), with six months to maturity and delivery price \(K=102\).
  Three months later, with three months left to maturity, the stock trades
  at \(S_t=105\). Find the current value of the long forward position.
\end{exparts}
\end{exercise}


\subsection*{Part II --- Completeness and the Second Fundamental Theorem}

\begin{exercise}[Medium --- Completing a four-state market may take two new assets]
One period, four states. A bond pays \(1\) in every state. A stock with
\(S_0=100\) pays \((70,90,110,130)\) across the four states. A call struck at
\(K=105\) pays \((0,0,5,25)\).
\begin{exparts}
  \item Explain, without any computation, why three traded assets can never
  span four states.
  \item Using the states with payoffs \((70,90,130)\) --- i.e., omitting the
  third state --- form the \(3\times3\) matrix of bond, stock, and call
  payoffs and compute its determinant. Conclude that the three assets are
  linearly independent.
  \item A second call, struck at \(K=115\), pays \((0,0,0,15)\). Form the
  \(4\times4\) payoff matrix of all four states and all four assets
  (bond, stock, call at \(105\), call at \(115\)), and compute its
  determinant.
  \item Conclude whether the market is complete once the second call is
  added, and relate your answer to the Second Fundamental Theorem.
\end{exparts}
\end{exercise}

\subsection*{Part III --- Superhedging and Price Bounds}

\begin{exercise}[Easy to medium --- Bounding a worst-of-two-stocks claim]
Let \(H=\min(S_T^1,S_T^2)\), the mirror image of the best-of-two-stocks claim
from lecture.
\begin{exparts}
  \item Show \(\min(a,b)\leq a\) and \(\min(a,b)\leq b\) for \(a,b\geq0\).
  Conclude that holding either stock alone \emph{superhedges} \(H\), and
  state the resulting upper bound on \(V_t(H)\).
  \item Explain why holding cash alone (\(Y_T\equiv0\)) is a valid but very
  loose subhedge. Why does a materially tighter lower bound generally
  require a claim on the \emph{spread} \(S_T^1-S_T^2\), rather than a
  position in the two stocks alone?
  \item \emph{Numerical.} With \(S_t^1=100\) and \(S_t^2=95\), state the
  resulting price corridor for \(H\).
\end{exparts}
\end{exercise}

